\documentclass{article}

\usepackage{amsmath,amssymb}
\usepackage[a4paper,margin=2.5cm]{geometry}

\title{Derivation of Governing Equations for Two-Dimensional Natural Convection}
\author{}
\date{}

\begin{document}

\maketitle

\section{Introduction}

In this study, a
two-dimensional natural convection model is considered for an
incompressible fluid. The governing equations are derived from the conservation of
mass, momentum, and thermal energy. The resulting equations
describe the fluid velocity, pressure, and temperature fields.

\section{Assumptions}

The following assumptions are adopted:

\begin{enumerate}
    \item The fluid is Newtonian and incompressible.
    \item The flow is laminar.
    \item The fluid properties are constant, except for density
    variation in the buoyancy term.
    \item The Boussinesq approximation is valid.
    \item The fluid contains no NEPCM particles and no phase-change
    effects are considered.
    \item There are no internal heat sources.
    \item The flow is two-dimensional, with velocity components
    $u$ and $v$ in the $x$ and $y$ directions, respectively.
\end{enumerate}

The coordinate system is defined such that the $x$-direction is
horizontal and the positive $y$-direction is vertically upward.
Gravity acts in the negative $y$-direction.

\section{Continuity Equation}

The continuity equation represents the conservation of mass.
The general mass conservation equation is

\begin{equation}
\frac{\partial \rho}{\partial t}
+\nabla\cdot(\rho\mathbf{u})=0,
\label{eq:mass_general}
\end{equation}

where $\rho$ is the fluid density, $t$ is time, and
$\mathbf{u}$ is the velocity vector.

For an incompressible fluid with constant reference density
$\rho_0$, the density is independent of time and position
in the continuity equation. And, $\rho_0$ = $\rho$ = constant. Therefore,

\begin{equation}
\frac{\partial \rho}{\partial t}=0,
\qquad
\nabla\cdot(\rho_0\mathbf{u})
=\rho_0\nabla\cdot\mathbf{u}.
\end{equation}

Substituting these expressions into
Equation~\eqref{eq:mass_general} gives

\begin{equation}
\rho_0\nabla\cdot\mathbf{u}=0.
\end{equation}

Since $\rho_0$ is non-zero, the continuity equation becomes

\begin{equation}
\nabla\cdot\mathbf{u}=0.
\end{equation}

For a two-dimensional velocity field
$\mathbf{u}=(u,v)$, the continuity equation is

\begin{equation}
\boxed{
\frac{\partial u}{\partial x}
+\frac{\partial v}{\partial y}=0
}
\label{eq:continuity}
\end{equation}

This equation ensures that the fluid satisfies mass
conservation without local expansion or compression.

\section{Momentum Equations}

The momentum equations are derived from Newton's second law,
which states that the rate of change of momentum is equal to
the net force acting on the fluid.

For a Newtonian fluid, the momentum conservation equation is

\begin{equation}
\rho\frac{D\mathbf{u}}{Dt}
=-\nabla p+\mu\nabla^2\mathbf{u}
+\rho\mathbf{g},
\label{eq:momentum_general}
\end{equation}

where $\rho$ is the fluid density, $p$ is pressure, $\mu$ is
dynamic viscosity, and $\mathbf{g}$ is the gravitational
acceleration vector.

The material derivative describes the acceleration of a
moving fluid particle and is defined as

\begin{equation}
\frac{D}{Dt}
=\frac{\partial}{\partial t}
+\mathbf{u}\cdot\nabla.
\end{equation}

For two-dimensional flow, the horizontal and vertical
acceleration components are

\begin{align}
\frac{Du}{Dt}
&=\frac{\partial u}{\partial t}
+u\frac{\partial u}{\partial x}
+v\frac{\partial u}{\partial y},
\\
\frac{Dv}{Dt}
&=\frac{\partial v}{\partial t}
+u\frac{\partial v}{\partial x}
+v\frac{\partial v}{\partial y}.
\end{align}

The momentum equations can be simplified by applying the
constant-viscosity and incompressibility assumptions. The
viscous terms become $\mu\nabla^2u$ and $\mu\nabla^2v$,
respectively.

\subsection{Boussinesq Approximation}

In natural convection, temperature differences cause small
variations in fluid density. Under the Boussinesq approximation,
density is approximated by

\begin{equation}
\rho\approx\rho_0
\left[1-\beta(T-T_{\mathrm{ref}})\right],
\label{eq:density}
\end{equation}

where $\rho_0$ is the reference density, $\beta$ is the
volumetric thermal expansion coefficient, $T$ is the local
temperature, and $T_{\mathrm{ref}}$ is the reference
temperature.

Density variations are neglected in the inertia and viscous
terms but retained in the gravitational force. After removing
the reference hydrostatic pressure contribution, the effective
buoyancy force per unit volume in the vertical direction is

\begin{equation}
f_b=\rho_0g\beta(T-T_{\mathrm{ref}}),
\label{eq:buoyancy}
\end{equation}

where $g$ is the magnitude of gravitational acceleration.
With the positive $y$-direction pointing upward, this term
is positive for fluid warmer than the reference temperature.

\subsection{Horizontal Momentum Equation}

In the horizontal direction, gravity has no direct component.
Therefore, the horizontal momentum equation is

\begin{equation}
\rho_0\frac{Du}{Dt}
=-\frac{\partial p}{\partial x}
+\mu\nabla^2u.
\end{equation}

Expanding the material derivative and the Laplacian gives

\begin{equation}
\boxed{
\begin{aligned}
\rho_0\left(
\frac{\partial u}{\partial t}
+u\frac{\partial u}{\partial x}
+v\frac{\partial u}{\partial y}
\right)
={}&-\frac{\partial p}{\partial x}\\
&+\mu\left(
\frac{\partial^2u}{\partial x^2}
+\frac{\partial^2u}{\partial y^2}
\right).
\end{aligned}
}
\label{eq:xmomentum}
\end{equation}

The pressure-gradient term represents the pressure force,
while the viscous term describes the diffusion of horizontal
momentum due to fluid viscosity.

\subsection{Vertical Momentum Equation}

In the vertical direction, the momentum equation includes
the buoyancy contribution derived using the Boussinesq
approximation. Thus,

\begin{equation}
\rho_0\frac{Dv}{Dt}
=-\frac{\partial p}{\partial y}
+\mu\nabla^2v
+\rho_0g\beta(T-T_{\mathrm{ref}}).
\end{equation}

Expanding the material derivative and the Laplacian gives

\begin{equation}
\boxed{
\begin{aligned}
\rho_0\left(
\frac{\partial v}{\partial t}
+u\frac{\partial v}{\partial x}
+v\frac{\partial v}{\partial y}
\right)
={}&-\frac{\partial p}{\partial y}\\
&+\mu\left(
\frac{\partial^2v}{\partial x^2}
+\frac{\partial^2v}{\partial y^2}
\right)\\
&+\rho_0g\beta(T-T_{\mathrm{ref}}).
\end{aligned}
}
\label{eq:ymomentum}
\end{equation}

The buoyancy term drives upward motion in regions warmer
than the reference temperature and downward motion in
regions colder than the reference temperature.

\section{Energy Equation}

The energy equation represents the conservation of thermal
energy. For a fluid with constant thermal conductivity and
specific heat capacity, with no internal heat generation,
the thermal energy equation is

\begin{equation}
\rho_0c_p\frac{DT}{Dt}
=k\nabla^2T,
\label{eq:energy_general}
\end{equation}

where $c_p$ is the specific heat capacity at constant
pressure, $k$ is thermal conductivity, and $T$ is temperature.

The material derivative of temperature is

\begin{equation}
\frac{DT}{Dt}
=\frac{\partial T}{\partial t}
+u\frac{\partial T}{\partial x}
+v\frac{\partial T}{\partial y}.
\end{equation}

Substituting this expression into
Equation~\eqref{eq:energy_general} gives

\begin{equation}
\rho_0c_p\left(
\frac{\partial T}{\partial t}
+u\frac{\partial T}{\partial x}
+v\frac{\partial T}{\partial y}
\right)
=k\nabla^2T.
\end{equation}

Dividing by $\rho_0c_p$ and defining the thermal
diffusivity as

\begin{equation}
\alpha=\frac{k}{\rho_0c_p},
\end{equation}

the energy equation becomes

\begin{equation}
\frac{\partial T}{\partial t}
+u\frac{\partial T}{\partial x}
+v\frac{\partial T}{\partial y}
=\alpha\nabla^2T.
\end{equation}

Expanding the Laplacian gives the final two-dimensional
energy equation:

\begin{equation}
\boxed{
\begin{aligned}
\frac{\partial T}{\partial t}
+u\frac{\partial T}{\partial x}
+v\frac{\partial T}{\partial y}
=\alpha\left(
\frac{\partial^2T}{\partial x^2}
+\frac{\partial^2T}{\partial y^2}
\right).
\end{aligned}
}
\label{eq:energy}
\end{equation}

The left-hand side represents the local and convective
changes in temperature, while the right-hand side represents
heat conduction within the fluid.

\section{Boundary Conditions}

For stationary, impermeable no-slip walls,
\begin{equation}
u=v=0
\quad \text{on all walls}.
\end{equation}

The hot and cold vertical walls satisfy
\begin{align}
T&=T_h &&\text{at }x=0,\\
T&=T_c &&\text{at }x=L.
\end{align}

The insulated horizontal walls satisfy
\begin{equation}
\frac{\partial T}{\partial y}=0
\quad\text{at }y=0,H.
\end{equation}

\end{document}
